\documentclass[12pt,a4paper]{article}

\usepackage[T2A]{fontenc}
\usepackage[utf8]{inputenc}
\usepackage[russian]{babel}
\usepackage{amsmath,amssymb,amsfonts}
\usepackage{geometry}
\usepackage{booktabs}
\usepackage{xcolor}

\geometry{left=22mm,right=18mm,top=18mm,bottom=20mm}
\setlength{\parindent}{1.25em}
\setlength{\parskip}{0.35em}

\newcommand{\Reynolds}{\operatorname{Re}}
\newcommand{\Laplace}{\Delta}
\newcommand{\J}{J}
\newcommand{\norm}[1]{\left\lVert #1 \right\rVert}
\newcommand{\inner}[2]{\left\langle #1,#2 \right\rangle}
\newcommand{\timing}[1]{\noindent\textcolor{gray}{\textbf{#1}}\quad}
\newcommand{\slide}[1]{%
    \par\noindent
    \fbox{\begin{minipage}{0.96\textwidth}
    \textbf{Слайд:} #1
    \end{minipage}}
    \par
}

\begin{document}

\begin{center}
    {\Large\bfseries Текст выступления на 10 минут}\\[0.5em]
    {\large Задача стабилизации для задачи магнитной гидродинамики\\
    в переменных \((\psi,\omega)\)}
\end{center}

\section*{0:00--0:30. Титульный слайд}

\slide{Название работы.}

\timing{0:00--0:30}
Рассматривается задача численного моделирования и стабилизации двумерного течения в переменных функция тока--завихренность.
Доклад состоит из двух частей.
Первая часть --- верификация численной реализации на задаче о квадратной каверне.
Вторая часть --- четырёхвихревая задача с внешним завихренным воздействием, поиск стационара, спектральный анализ и удаление неустойчивой компоненты.

\section*{0:30--1:30. Постановка задачи}

\slide{Уравнения Навье--Стокса и переход к переменным \((\psi,\omega)\).}

\timing{0:30--1:30}
Исходная система --- уравнения Навье--Стокса для несжимаемой вязкой жидкости:
\[
    \frac{\partial \mathbf u}{\partial t}
    +
    (\mathbf u\cdot\nabla)\mathbf u
    =
    -\nabla p+\nu\Laplace\mathbf u+\mathbf f,
    \qquad
    \nabla\cdot\mathbf u=0.
\]
Здесь \(\mathbf u=(u,v)\), \(p\) --- давление, \(\nu=1/\Reynolds\).
Чтобы исключить давление, выполняется переход к функции тока и завихренности.
Скорость задаётся через функцию тока:
\[
    u=\psi_y,\qquad v=-\psi_x.
\]
Тогда условие \(\operatorname{div}\mathbf u=0\) выполняется автоматически.
Завихренность связана с функцией тока уравнением
\[
    \omega=v_x-u_y,\qquad \Laplace\psi+\omega=0.
\]
После применения операции ротора к уравнениям Навье--Стокса получается уравнение для завихренности:
\[
    \omega_t
    =
    \nu\Laplace\omega
    -
    \J(\psi,\omega)
    -
    F_\omega(x,y),
\]
где
\[
    \J(\psi,\omega)=\psi_x\omega_y-\psi_y\omega_x.
\]
Итак, численно решается система
\[
    \begin{cases}
    \Laplace\psi+\omega=0,\\
    \omega_t=\nu\Laplace\omega-\J(\psi,\omega)-F_\omega(x,y).
    \end{cases}
\]
В задаче о каверне \(F_\omega=0\), а в четырёхвихревой задаче эта правая часть задаётся синусоидальными модами.

\section*{1:30--2:15. Граничные условия и якобиан}

\slide{Общие граничные условия. Откуда возникает якобиан.}

\timing{1:30--2:15}
На твёрдой непроницаемой стенке функция тока постоянна.
Для рассматриваемых замкнутых областей константа выбирается равной нулю:
\[
    \psi|_{\partial\Omega}=0.
\]
Завихренность на стенке не задаётся независимо.
Она восстанавливается по формуле Тома.
Для стенки с тангенциальной скоростью \(U\)
\[
    \omega_{\mathrm{wall}}
    =
    -\frac{2\psi_{\mathrm{near}}}{h^2}
    -
    \frac{2U}{h}.
\]
Частные граничные условия для каждой задачи задаются отдельно: в каверне движется верхняя стенка, а в четырёхвихревом эксперименте все стенки неподвижны.

Нелинейный член возникает из переноса завихренности скоростью:
\[
    (\mathbf u\cdot\nabla)\omega=u\omega_x+v\omega_y.
\]
После подстановки \(u=\psi_y,\ v=-\psi_x\) он записывается через якобиан \(\J(\psi,\omega)\).
В задаче о каверне этот член аппроксимируется центральными разностями, а во второй задаче заменяется схемой Аракавы.

\section*{2:15--3:30. Разностная схема}

\slide{Разностные операторы, факторизация и прогонка.}

\timing{2:15--3:30}
На равномерной сетке используются центральные разности:
\[
    (D_x f)_{ij}=\frac{f_{i+1,j}-f_{i-1,j}}{2h_x},
    \qquad
    (D_y f)_{ij}=\frac{f_{i,j+1}-f_{i,j-1}}{2h_y}.
\]
Лапласиан аппроксимируется пятиточечным оператором.

Для уравнения на функцию тока вводится псевдовременной процесс
\[
    \psi_t=\Laplace_h\psi+\omega.
\]
Неявный шаг факторизуется:
\[
    (I-\tau D_{xx})(I-\tau D_{yy})\psi^{n+1}
    =
    \psi^n+\tau\omega^n+\tau^2D_{xx}D_{yy}\psi^n.
\]
Здесь важен множитель
\[
    (I-\tau D_{xx})(I-\tau D_{yy})
    =
    I-\tau\Laplace_h+\tau^2D_{xx}D_{yy}.
\]
Дополнительный член с \(\tau^2D_{xx}D_{yy}\) компенсирует смешанный член факторизации в стационарном пределе.

Для завихренности используется аналогичная факторизация:
\[
    (I-\tfrac{\tau}{\Reynolds}D_{xx})
    (I-\tfrac{\tau}{\Reynolds}D_{yy})\omega^{n+1}
    =
    \omega^n-\tau\J_h-\tau F_\omega
    +
    \frac{\tau^2}{\Reynolds^2}D_{xx}D_{yy}\omega^n.
\]
После факторизации двумерная задача распадается на последовательность трёхдиагональных систем, которые решаются методом прогонки.
Для функции тока это две системы:
\[
    (I-\tau D_{xx})\psi^*
    =
    \psi^n+\tau\omega^n+\tau^2D_{xx}D_{yy}\psi^n,
    \qquad
    (I-\tau D_{yy})\psi^{n+1}=\psi^*.
\]
Для завихренности аналогично:
\[
    (I-\tfrac{\tau}{\Reynolds}D_{xx})\omega^*
    =
    \omega^n-\tau\J_h-\tau F_\omega
    +
    \frac{\tau^2}{\Reynolds^2}D_{xx}D_{yy}\omega^n,
    \qquad
    (I-\tfrac{\tau}{\Reynolds}D_{yy})\omega^{n+1}=\omega^*.
\]
Вычислительный шаг устроен так: сначала находится \(\psi\), затем восстанавливается граничная завихренность по формуле Тома, затем обновляется \(\omega\), после чего вычисляются скорости и стационарные невязки.

\section*{3:30--4:00. Схема Аракавы}

\slide{Девятиточечная аппроксимация якобиана.}

\timing{3:30--4:00}
Для четырёхвихревой задачи используется схема Аракавы:
\[
    \J_h^A(\psi,\omega)
    =
    \frac{1}{3}
    \left(
    \J_h^{(1)}+\J_h^{(2)}+\J_h^{(3)}
    \right).
\]
Это усреднение трёх разностных форм одного и того же якобиана.
Оно сохраняет второй порядок по пространству, использует девятиточечный шаблон и уменьшает нефизическую перестройку вихревой картины на длинных расчётах.

\section*{4:00--5:20. Каверна и сравнение с литературными данными}

\slide{Постановка каверны, линии тока, три слайда сравнения центров.}

\timing{4:00--5:20}
Верификационная задача --- квадратная каверна:
\[
    \Omega=(0,1)\times(0,1),\qquad F_\omega=0.
\]
На границе
\[
    \psi|_{\partial\Omega}=0,
\]
верхняя стенка движется со скоростью \(U_{\mathrm{top}}=1\), остальные стенки неподвижны.
Расчёты выполнены на сетке \(601\times601\) для трёх чисел Рейнольдса:
\[
    \Reynolds=1000,\quad 5000,\quad 10000.
\]
Сравнение проводится по координатам центров вихрей и значениям функции тока.

\begin{center}
\begin{tabular}{c c c c}
\toprule
\(\Reynolds\) & центр главного вихря & центр вторичного вихря & \(\psi_{\min}\)\\
\midrule
1000 & \((0.5300;\ 0.5650)\) & \((0.8633;\ 0.1117)\) & \(-0.118781\)\\
5000 & \((0.5150;\ 0.5350)\) & \((0.8067;\ 0.0733)\) & \(-0.118382\)\\
10000 & \((0.5117;\ 0.5300)\) & \((0.7767;\ 0.0583)\) & \(-0.098797\)\\
\bottomrule
\end{tabular}
\end{center}

В презентации для каждого числа Рейнольдса оставлено место под график из статьи.
Под графиками приведены численные значения расчёта и значения Erturk et al.
На профилях скоростей квадратными маркерами дополнительно показаны табличные точки Ghia et al. для центральных сечений \(u(0.5,y)\) и \(v(x,0.5)\).

\section*{5:20--5:55. Стационарная задача и переход ко второй части}

\slide{Стационарная система и задача стабилизации.}

\timing{5:20--5:55}
После проверки схемы на каверне рассматривается стационарная задача
\[
    \Laplace\psi_*+\omega_*=0,
\]
\[
    \nu\Laplace\omega_*-\J(\psi_*,\omega_*)-F_\omega(x,y)=0.
\]
Если около стационара есть собственное значение линеаризованного оператора с положительной действительной частью, стационар является неустойчивым.
Поэтому задача стабилизации формулируется как выделение и удаление растущей компоненты возмущения.

\section*{5:55--6:45. Четырёхвихревой эксперимент}

\slide{Постановка и выбор начального приближения.}

\timing{5:55--6:45}
Вторая задача решается в прямоугольнике
\[
    \Omega=(0,1)\times(0,0.5),\qquad \Reynolds=260.
\]
Частные граничные условия:
\[
    \psi|_{\partial\Omega}=0,\qquad U_{\partial\Omega}=0.
\]
Четырёхвихревая структура создаётся внешней силой
\[
    F_\omega(x,y)
    =
    \sum_{(m,n)\in\mathcal M}
    a_{mn}\frac{\pi n}{L_y}
    \sin\frac{\pi m x}{L_x}
    \sin\frac{\pi n y}{L_y}.
\]
Прямое интегрирование из нулевого состояния даёт последовательность сохранённых снимков.
Кандидат в стационар выбирается как снимок, на котором достигается минимум разности по сохранённым результатам.
В расчёте это шаг \(236000\), время \(t=2.36001\).

\section*{6:45--7:35. Метод Ньютона}

\slide{Ньютон для нелинейной стационарной системы.}

\timing{6:45--7:35}
Кандидат из прямого интегрирования является только приближением к стационару.
Далее решается нелинейная алгебраическая система
\[
    F(U)=0,\qquad U=(\psi,\omega)^T.
\]
Вектор \(F(U)\) состоит из двух блоков: невязки уравнения Пуассона и невязки стационарного уравнения переноса завихренности во внутренних узлах сетки.

На \(k\)-й итерации метода Ньютона линеаризуется \(F\):
\[
    F(U_k+\delta U_k)
    \approx
    F(U_k)+F'(U_k)\delta U_k.
\]
Поправка находится из системы
\[
    F'(U_k)\delta U_k=-F(U_k),
\]
после чего выполняется обновление
\[
    U_{k+1}=U_k+\alpha_k\delta U_k.
\]
Коэффициент \(\alpha_k\) выбирается поиском по шагу, чтобы уменьшить норму невязки.
Линейная система решается методом GMRES, причём матрица Якоби явно не собирается: используется только действие \(F'(U_k)\) на вектор.

\section*{7:35--8:35. Линеаризация и поиск собственных чисел}

\slide{Вывод оператора и матрично-свободный спектральный расчёт.}

\timing{7:35--8:35}
Пусть найденный стационар равен \(U_*=(\psi_*,\omega_*)\).
Он удовлетворяет стационарной системе
\[
    \Laplace\psi_*+\omega_*=0,
    \qquad
    \nu\Laplace\omega_*-\J(\psi_*,\omega_*)-F_\omega=0.
\]
Рассматривается малое возмущение
\[
    \psi=\psi_*+\phi,\qquad
    \omega=\omega_*+\eta.
\]
Линеаризация делается напрямую: подставляем \(\psi_*+\phi\), \(\omega_*+\eta\) в исходную систему,
вычитаем уравнения для стационара и оставляем только линейные по \(\phi,\eta\) члены.
Нелинейный якобиан раскладывается как
\[
    \J(\psi_*+\phi,\omega_*+\eta)
    =
    \J(\psi_*,\omega_*)
    +
    \J(\phi,\omega_*)
    +
    \J(\psi_*,\eta)
    +
    \J(\phi,\eta).
\]
Первый член сокращается со стационарным уравнением, последний является квадратично малым.
Поэтому получается
\[
    \Laplace\phi+\eta=0,
\]
\[
    \eta_t
    =
    \nu\Laplace\eta
    -
    \J(\phi,\omega_*)
    -
    \J(\psi_*,\eta).
\]
Из первого уравнения, при однородном граничном условии для возмущения функции тока,
\[
    \phi=-\Laplace^{-1}\eta.
\]
Поэтому линейная динамика записывается как
\[
    \eta_t=\mathcal A\eta,
\]
\[
    \mathcal A\eta=
    \nu\Laplace\eta
    -
    \J(-\Laplace^{-1}\eta,\omega_*)
    -
    \J(\psi_*,\eta).
\]
Собственные числа ищутся из задачи
\[
    \mathcal A\varphi=\lambda\varphi.
\]
Здесь используется подстановка нормальной моды
\(\eta(x,y,t)=e^{\lambda t}\varphi(x,y)\):
производная по времени превращается в умножение на \(\lambda\), а общий множитель \(e^{\lambda t}\) сокращается.
Матрица \(\mathcal A\) явно не собирается.
Алгоритм получает только действие \(\eta\mapsto\mathcal A\eta\): для этого решается вспомогательная задача Пуассона для \(\phi\), затем вычисляются два линеаризованных якобиана.
Собственные значения с наибольшей действительной частью находятся методом Арнольди.

\section*{8:35--9:20. Неустойчивая мода и коррекция}

\slide{Спектр, мода, вычитание проекции.}

\timing{8:35--9:20}
В спектре найдена одна комплексно-сопряжённая неустойчивая пара:
\[
    \lambda_{0,1}=0.217011\pm1.786819\,i.
\]
Положительная действительная часть означает линейную неустойчивость.
Мнимая часть задаёт частоту колебаний.

Если комплексная мода имеет вид \(v=a+ib\), то вещественное неустойчивое подпространство:
\[
    E_{\mathrm{unst}}
    =
    \operatorname{span}\{\operatorname{Re}v,\operatorname{Im}v\}.
\]
Для отклонения от стационара
\[
    r(t)=U(t)-U_*
\]
вычисляется проекция
\[
    r_{\mathrm{unst}}
    =
    \sum_{\ell=1}^{m}
    \inner{r(t)}{e_\ell}e_\ell.
\]
Скорректированное состояние:
\[
    U_{\mathrm{corr}}(t)=U(t)-r_{\mathrm{unst}}(t).
\]
Дополнительно сравниваются координаты центров четырёх вихрей.
В прямом расчёте центры диагональных ячеек смещаются к центральной линии, а после коррекции возвращаются ближе к кандидату в стационарное состояние.
Также построен график относительного отклонения всего состояния \(U=(\psi,\omega)\) от четырёхвихревого кандидата
\[
    U_{\mathrm{cand}}=U(t_{\mathrm{cand}}),
    \qquad
    t_{\mathrm{cand}}=2.36001,\quad n_{\mathrm{cand}}=236000.
\]
Этот кандидат выбран как сохранённый снимок с минимальной разностью между соседними сохранёнными состояниями.
На графике показана величина
\[
    d(t)=
    \frac{\norm{U(t)-U_{\mathrm{cand}}}_2}
         {\norm{U_{\mathrm{cand}}}_2}.
\]
Эта величина лучше отражает эффект коррекции, чем скорость движения центров, потому что учитывает всё поле, а не только дискретно выбранные экстремумы.

\section*{9:20--10:00. Итоги}

\slide{Итоговый слайд.}

\timing{9:20--10:00}
Построена численная схема в переменных \((\psi,\omega)\) с факторизацией по направлениям и прогонками.
Для задачи о каверне на сетке \(601\times601\) получено совпадение координат центров вихрей с литературными данными.
Для четырёхвихревой задачи найден стационар методом Ньютона и построен линеаризованный оператор.
Спектральный расчёт выявил неустойчивую пару
\[
    \lambda=0.217011\pm1.786819i.
\]
Вычитание проекции на соответствующее неустойчивое подпространство уменьшает уход решения от четырёхвихревой структуры.

\end{document}
